\documentclass[10pt,letterpaper]{article}
\usepackage[letterpaper,margin=1.6cm]{geometry}
\usepackage[T1]{fontenc}
\usepackage[utf8]{inputenc}
\usepackage{newtxtext}
\usepackage{xcolor}
\usepackage[hidelinks]{hyperref}
\usepackage{enumitem}
\pagestyle{empty}
\setlength{\parindent}{0pt}
\setlength{\parskip}{2pt}
\linespread{0.96}
\setlist[itemize]{leftmargin=5mm,itemsep=0pt,topsep=1pt,parsep=0pt}
\definecolor{ieeeblue}{RGB}{0,84,166}
\newcommand{\clhead}[1]{\textcolor{ieeeblue}{\textbf{#1}}}

\begin{document}

\begin{flushright}
\textbf{St\'ephane Ga\"el R. Ekodeck}\\
Department of Computer Science, LIRIMA, Team GRIMCAPE\\
University of Yaound\'e I, P.O.\ Box 812, Yaound\'e, Cameroon\\
\href{mailto:stephane-gael.ekodeck@facsciences-uy1.cm}%
{stephane-gael.ekodeck@facsciences-uy1.cm}\\
June 2026
\end{flushright}

The Editor-in-Chief\\
\textit{IEEE Transactions on Information Forensics and Security}\\
IEEE Signal Processing Society\\

Dear Editor-in-Chief,

We respectfully submit our manuscript, \textbf{``NARCIS: Authenticated
Neural Coverless Image Signaling with Attack-Qualified Codebooks and
Error Correction,''} for consideration as an original research article in
\textit{IEEE Transactions on Information Forensics and Security} (TIFS).

\clhead{Scope alignment.}
NARCIS is an end-to-end security protocol that transmits authenticated
messages by selecting unchanged natural images from a shared index.
Its core contributions are at the intersection of information forensics
and security: authenticated encryption of payload and metadata, formal
security propositions grounded in AES-GCM's established properties,
selection-bias measurement with three detector classes, and systematic
evaluation of an attack-qualified finite-index communication system.
These contributions fit squarely within TIFS's focus on cryptographic
applications, multimedia security, and forensics.

\clhead{Key technical contributions.}
\begin{itemize}
\item \textbf{Feasible operating point (Contribution 1).} A codebook is
  accepted only when every attack-qualified bucket can satisfy the
  multiplicity required by every protected coded message. The selected
  alphabet is the largest satisfying this criterion. Over 1{,}575
  end-to-end message-condition trials on BOSSBase~1.01 and Caltech-101,
  NARCIS recovered 100\% of authenticated plaintexts under 21 channel
  conditions including eight unseen holdout strengths.
\item \textbf{End-to-end authenticated protocol (Contribution 2).}
  AES-GCM~(NIST SP 800-38D) encrypts the payload; associated data binds
  sender and receiver identities. A monotonic sequence counter rejects
  replay. A keyed Gray permutation limits public symbol-to-label
  correspondence. Reed--Solomon forward error correction handles residual
  label errors. Two formal propositions (Confidentiality;
  Authenticity+Freshness) ground these properties in AES-GCM's
  IND-CCA2 and authentication guarantee. Net throughput at a 16-symbol
  operating point scales from 0.184 to 1.213 bits per cover for
  8- to 128-byte payloads.
\item \textbf{Systematic detectability and ablation (Contribution 3).}
  Three detector classes (global logistic regression, residual-feature
  ExtraTrees, selection CNN) are evaluated on five partitions of each
  dataset. BOSSBase AUCs include chance. Caltech-101 shows a weak but
  reproducible partition-specific CNN signal (20-repeat mean AUC~0.533,
  95\%~CI [0.517, 0.548]), reported transparently. Removing attack
  qualification reduces BOSSBase seed-11 complete-message success from
  210/210 to 142/210~(67.6\%), with crop and large-rotation conditions
  failing completely, demonstrating that the qualification filter is the
  dominant robustness mechanism.
\end{itemize}

\clhead{Novelty and positioning.}
Selection-based coverless systems have been evaluated at the
descriptor level (average symbol accuracy, nominal capacity) but have
not integrated cryptographic payload protection, replay control, or
error correction into a single verifiable system. NARCIS introduces the
\emph{feasible-codebook} criterion as the primary evaluation unit, closes
this protocol gap, and provides matching detectability and ablation
evidence. The paper explicitly positions itself relative to related
methods (Abdulsattar~2021; Guo--Ping~2026; Ren--Wu~2024) without
overstating numerical comparisons across different datasets and protocols.

\clhead{Author declarations.}
\begin{itemize}
\item The manuscript is original, has not been published, and is not under
  consideration elsewhere. It was previously desk-rejected by JISA
  (Ms.\ No.\ JISAS-D-26-03622) for scope reasons; this submission has
  been substantially revised and is not a transfer.
\item All four authors approved the final manuscript and author order.
  WYSAWIS (cited in related work) shares authors with this submission
  and is identified explicitly in the paper; it is used as qualitative
  context, not as an experimental baseline.
\item No competing interests and no external funding.
\item No human subjects, personal data, or animals are involved.
  Implementation, fixed seeds, experimental results, and manuscript
  sources are available at
  \href{https://github.com/EkodeckStephane/NARCIS}%
  {github.com/EkodeckStephane/NARCIS}.
\end{itemize}

We believe NARCIS will interest TIFS readers working on authenticated
multimedia communication, selection forensics, and practical steganographic
security systems. We welcome the opportunity to provide any additional
information the Associate Editor may require.\\[4pt]

Yours sincerely,\\[6pt]

\textbf{St\'ephane Ga\"el R. Ekodeck} (Corresponding Author)\\
ORCID: 0000-0002-8094-8832\\
\textbf{Co-authors:} Vivien Beyala Kamgang (ORCID 0000-0002-3644-1866);
Serge Alain Eb\'el\'e (ORCID 0000-0001-7036-7068);
Julius Marcellin Nkenlifack (ORCID 0000-0003-3322-4687).

\end{document}
